Selective state space models (SSMs) offer global context at linear cost, but their cost inside lightweight detectors on embedded CPUs is poorly characterized. We replace the $P_5$ self-attention block (\texttt{C2PSA}) of YOLO11s with a plugin-free, four-direction 2D selective scan (SS2D) built from standard operators, and study its accuracy on a cross-facility partition of the synthetic NVIDIA PhysicalAI-SmartSpaces warehouse benchmark and its deployment cost on a Raspberry Pi~4. The replacement removes $6.72\%$ of the parameters and $0.5$ GFLOPs, mainly by narrowing channels. In preliminary single-seed experiments with single-stage, two-stage, and dual-scale training, it does not preserve accuracy: the baseline leads at the final epoch in both $\text{mAP}_{50}$ and $\text{mAP}_{50-95}$. We call a model edge-feasible if it runs in FP32 ONNX Runtime without custom operators within $1$\,GiB of memory; the SS2D model meets this criterion but is $18\%$ slower than the baseline with four threads ($1314$ vs.\ $1117$\,ms forward latency) and uses $60\%$ more peak memory despite fewer FLOPs. Operator profiling and a runtime intervention show that the overhead comes from how the scan is executed rather than from its arithmetic or graph size: re-expressing the same scan as an ONNX \texttt{Loop} with $40\times$ fewer nodes, as a chunked scan, or as a parallel associative scan without the sequential chain made the model slower, not faster. We summarize design lessons for hybrid CNN--SSM detectors on edge CPUs, including a pitfall in zero-initializing a convolution ahead of batch normalization.
